\documentclass[../main.tex]{subfiles}
\graphicspath{{\subfix{../images/}}}
\begin{document}

\chapter{Introduction}

Facial expression recognition (FER) plays a crucial role within human-computer interaction (HCI), enabling automated detection of human emotions from facial data \cite{li2022survey}.
FER has also attracted researchers in the field of affective computing, due to its potential applications in creating systems that detect, understand, and replicate emotions \cite{karnati2023survey}.
Thanks to advancements in computer vision and deep learning (DL) methods, FER technology has improved significantly over the recent years \cite{li2022survey}.
This thesis focuses on the subfield of lightweight FER for deployment on resource-constrained environments, which demands a balance between accuracy and efficiency.\\

FER systems have a wide range of applications, including sociable robots, medical treatment, driver safety, marketing, education, psychological research, and security \cite{li2022survey, canal2022survey, karnati2023survey, sajjad2023survey}.
Despite this, FER faces several challenges such as occlusion, pose and lighting conditions, subject variability, data annotation, resolution, and more \cite{li2022survey, sajjad2023survey}.
On top of this, real-world usage scenarios such as edge computing require the highest level of efficiency in order to be feasible \cite{wu2020edge}.\\

Given the proliferation of mobile and edge platforms with a limited resource budget, there is an increasing need for lightweight FER models.
Recent studies have focused on creating optimized deep learning methods for computer vision \cite{howard2017mobilenet, zhang2018shufflenet}, with some even applying such methods directly to FER \cite{eff-face2021, lengwe2023pattlite}.
These studies have shown that it is possible to significantly optimize these models while maintaining or even improving the system's accuracy.\\

This thesis proposal aims to enhance the lightweight FER network, EfficientFace \cite{eff-face2021}, in order to improve its accuracy while maintaining its low resource cost.
Part of this process involves modifying the network's training method by employing a high-cost, high-accuracy FER model \cite{apvit2023} that can be discarded during inference.
The rest of this document lays out the context and planification of the associated scientific study.\\

We now describe the structure and remaining chapters of this thesis proposal. 
Chapter 2 serves as the theoretical framework on which our proposal is based.
Chapter 3 formally presents the problem and proposed approach to treat it.
Chapter 4 defines the scope of the associated research project.
Lastly, chapter 5 provides a detailed planification of the study intended to test the validity of our proposal.

\end{document}